\NumberedItem{4}
\noindent{\small[S. 291]}\par
\noindent\hspace*{\parindent}Setzt man
\[
\begin{aligned}
&1+\frac{\alpha\beta}{\gamma\delta}x+\frac{\alpha{.}\,\alpha+1{.}\,\beta{.}\,\beta+1}{\gamma{.}\,\gamma+1{.}\,\delta{.}\,\delta+1}xx+\text{u.s.w.}\\[.4em]
&\qquad=\psi(\alpha,\beta,\gamma,\delta,x)=P,
\end{aligned}
\]
\noindent so ist
\[
\begin{aligned}
&(\alpha\beta-(\gamma-1)(\delta-1))P+((\alpha+\beta+1)x-(\gamma+\delta-1))\frac{dP}{dx}\\[.5em]
&\qquad+(xx-x)\frac{ddP}{dx^2}\\[.8em]
&=\frac{(\alpha\beta-\gamma\delta)(\gamma-1)(\delta-1)}{\gamma\delta}\\[.5em]
&\qquad+\frac{\alpha\beta(\gamma-1)(\delta-1)((\alpha+1)(\beta+1)-(\gamma+1)\delta+1))}{\gamma{.}\,\gamma+1{.}\,\delta{.}\,\delta+1}x\\[.8em]
&=A+Bx.
\end{aligned}
\]
\noindent Schreibt man
\[
x=\frac{t}{t-1},
\]
\noindent so wird die Gleichung
\[
\begin{aligned}
&\frac{\alpha\beta-(\gamma-1)(\delta-1)}{t-1}P-((\alpha+\beta-\gamma-\delta+2)t+\gamma+\delta-1)\frac{dP}{dt}\\[.5em]
&\qquad+(tt-t)\frac{ddP}{dt^2}\\[.8em]
&=\frac{-A+(A+B)t}{(t-1)^2}\mathpunct{\raisebox{.35ex}{.}}
\end{aligned}
\]
\noindent Diese Art führt nicht bequem zum Ziel. Wohl aber die auf S. 290 [des Handbuchs, siehe art. 3] angedeutete.
\clearpage
